\chapter{Anexo 3. Estimación de costos marginales implícitos con variables climáticas alternativas}
\label{chap:anexo3}

Esta sección explora la estimación de los costos marginales implícitos del sistema utilizando una construcción alternativa de las variables de estrés térmico, a fin de revisar la sensibilidad de la estimación ante cambios en las variables explicativas.

En particular, para esta estimación se sigue el mismo procedimiento señalado en la sección N°\ref{chap:datos}, con dos salvedades. Primero, se considera un umbral de 10°C para la exposición al frío, en lugar de 15°C en la estimación original. La razón detrás de esta modificación guarda relación con la posibilidad de que la variable esté relacionada con la exposición a temperaturas óptimas, lo que explicaría el la dirección del coeficiente.

Segundo, en lugar de considerar el promedio ponderado por hectáreas de cultivo, se considera el promedio simple para cada país. Esta modificación busca instalar robustez ante el cambio que el uso del suelo pudo haber tenido en el tiempo. Debido a que las hectáreas de cultivo de cada especie corresponden al año 2020, bien pueden ser reflejo de un proceso de adaptación en las décadas anteriores.

La Tabla N°\ref{tab:estimacion_costos_2} presenta los resultados tras estimar los costos marginales implícitos del sistema con el nuevo conjunto de variables de estrés térmico. Se observan resultados consistentes con la estimación predilecta, sobre todo para las variables de exposición al calor. Por otro lado, tanto la dirección de los coeficientes como la significancia estadística de las variables de exposición al frío impiden que pueda ser considerada como candidata para la estimación de escenarios contrafactuales.

\begin{table}[H]
   \caption{Estimación de costos marginales}
   \label{tab:estimacion_costos_2}
   \bigskip
   \centering
   \begin{tabular}{lccc}

      \hline 
      
      Var. Dep.: Costos marginales implícitos                     & (1) & (2) & (3)\\   
      \hline

      Tendencia             & -9,23$^{***}$ & -6,86$^{***}$ & -9,18 $^{***}$ \\   
                           & (0,8735) & (0,6910) & (0,8700) \\   
      Fertilizantes        & 0,49$^{***}$ & 0,48$^{***}$ & 0,49$^{***}$ \\   
                           & (0,0193) & (0,0197) & (0,0195)\\   
      Exposición al calor (t)      & 0,68$^{***}$ & - & 0,65$^{***}$ \\   
                           & (0,1541) & - & (0,1518) \\   
      Exposición al calor (t-1)      & 0,29$^{.}$ & - & 0,31$^{.}$ \\   
                           & (0,1622) & - & (0,1658)\\   
      Exposición al frío (t)      & - & -3,65 & -2,98$^{***}$ \\   
                           & - & (3,8850) & (3,5735)\\   
      Exposición al frío (t-1)      & - & 5,15 & 4,61\\   
                           & - & (3,2780) & (3,5061)\\  
      \hline
      Efecto fijo por país & $\checkmark$ & $\checkmark$ & $\checkmark$ \\  
      \hline
      Observaciones          & 1.282 & 1.282 & 1.282\\  
      R$^2$                 & 0,2800 & 0,2715 & 0,2797\\  
      Within R$^2$          & 0,1338 & 0,1238 & 0,1351\\        
      \bottomrule
   \end{tabular}
   \par \raggedright 
\end{table}

